\section{Introduction}

Machine-learning models rank every stock. Investors should not trade every rank difference.

A return-prediction model maps the information available at date $t$ into forecasts
\[
\widehat\mu_{1,t},\ldots,\widehat\mu_{N_t,t}.
\]
Portfolio formation then converts these forecasts into a total order, even when two predicted returns differ by only a few basis points. The ranking has no ``I do not know'' state. A stock must be placed above or below every other stock, and a rank-based strategy can turn each mechanical distinction into capital and turnover.

This paper asks when a predicted cross-sectional ordering is reliable enough to trade. The object of interest is not forecast accuracy alone and is not a confidence interval around each individual expected return. It is the risk of the investment decision induced by a rank relation. We distinguish
\[
\text{predicted ranking}
\qquad\text{from}\qquad
\text{reliable ranking}.
\]
The first is necessarily complete. The second can be a partial order: for some pairs the evidence supports $i\succ j$, for others it supports $j\succ i$, and for the remaining pairs the investor abstains.

This distinction matters because machine learning has made high-dimensional cross-sectional return prediction a canonical empirical-asset-pricing problem. \citet{GuKellyXiu2020} compare a broad set of linear and nonlinear methods using a large stock-characteristic panel and show the importance of nonlinear interactions. The subsequent literature increasingly evaluates the investment consequences of machine-learned forecasts. \citet{Allena2026} derives ex-ante confidence intervals for risk-premium forecasts and exploits cross-sectional differences in forecast precision, while \citet{JensenKellyMalamudPedersen2026} studies portfolio learning on an implementable, net-of-cost frontier. These developments raise the bar. ``Machine learning plus uncertainty'' is not by itself a contribution. Our object is different: forecast precision is an input, while ranking reliability is the expected loss of the pairwise decisions the investor actually deploys.

We introduce a nested family of partial rankings. Let $\widehat s_{i,t}>0$ be a predetermined scale summarizing the local difficulty of using stock $i$'s forecast. At abstention threshold $q\geq0$, stock $i$ is declared superior to stock $j$ only when
\[
\widehat\mu_{i,t}-q\widehat s_{i,t}
>
\widehat\mu_{j,t}+q\widehat s_{j,t}.
\]
At $q=0$, absent ties, the relation coincides with the model's total order. As $q$ increases, small or locally noisy forecast differences are removed. The surviving relation is a strict partial order, and the family is nested.

This construction gives a new measure of the investment opportunity set. We define \emph{reliable breadth} as the fraction of economically relevant pairwise rank relations retained at the selected threshold. Nominal breadth counts every distinction generated by the model. Reliable breadth counts only distinctions that the investor has enough evidence to monetize under a prespecified error budget. The time series of reliable breadth can therefore reveal when the cross-section is easy or difficult to rank and where the reliable information is located.

The statistical problem is unusual. For each candidate threshold, the investor observes one monthly sequence of cross-sectional ranking losses. These losses inherit persistence from expected returns, volatility, liquidity, common shocks, and the forecasting algorithm itself. Treating the months as independent can understate uncertainty precisely when the investment environment is most persistent. At the same time, estimating a mixing coefficient or a long-run covariance separately for every candidate threshold is unstable and operationally unattractive.

We use the risk-certification framework developed in the companion paper \citet{Bignozzi2026SNRCPS}. The natural selected-pair error rate is not automatically monotone as abstention increases: removing easy and difficult pairs can change its denominator in either direction. We therefore certify its smallest monotone majorant. For candidate thresholds $q_1<\cdots<q_K$, let $e_t(q_k)$ be the error rate among active relations in month $t$, and define
\[
L_t(q_k)
=
\max_{\ell\geq k}e_t(q_\ell).
\]
Then $L_t(q_k)$ is bounded, nonincreasing in $k$, and dominates the economically transparent loss $e_t(q_k)$. Self-normalization uses the complete certification path to select the least-abstaining threshold whose future risk is below $\alpha$ with confidence $1-\delta$. It requires neither a block length nor a long-run variance estimate.

A practical finance implementation cannot freeze one neural network for decades. We instead certify a predetermined prequential learning rule. Each forecast is generated using a fixed-length trailing window and no future data. The complete monthly loss is therefore a measurable finite-memory function of the underlying financial process. Under the sequential limit condition stated below, the certificate controls the future decision risk of the algorithm--threshold pair, while the model can continue to update according to its predeclared rule.

We then connect abstention to portfolio choice. Each certified relation can be represented as a relative-value edge. Aggregating edge out-degree minus in-degree gives a dollar-neutral portfolio. A fixed-denominator representation yields an exact decomposition of the raw machine-learned portfolio into certified and removed edges before leverage normalization. In a simple robust relative-value problem, the optimal position is zero whenever the lower end of the predicted spread fails to cover proportional trading costs. Thus, a partial ranking is not a statistical cosmetic: it is the discrete counterpart of an economically optimal no-trade region.

The empirical analysis is organized around four questions.

\begin{enumerate}[leftmargin=2em]
\item \textbf{How much?} What share of the total machine-learned ranking is reliably distinguishable, including within the economically relevant tails?
\item \textbf{When?} How does reliable breadth vary with aggregate volatility, market liquidity, funding conditions, and common-return variation?
\item \textbf{Where?} Which stocks, characteristics, sectors, and forecast disagreements account for reliable and unreliable relations?
\item \textbf{Why does it matter?} How much gross return, turnover, cost, capacity, and net economic value comes from certified versus uncertified rank relations?
\end{enumerate}

The baseline data are the preconstructed US stock-month panel of \citet{JensenKellyPedersen2023} distributed through WRDS, using the published 153-characteristic taxonomy and its already aligned next-month excess return. CRSP daily price and volume supplement the panel for capacity-aware cost measurement. Forecasts are generated recursively by ridge, elastic net, tree, boosting, neural-network, and ensemble models. The headline portfolio evidence is net of linear and capacity-aware trading costs, and every result is reported both with endogenous exposure reduction and after equal-gross rescaling.

The paper contributes at three levels. First, it introduces reliable breadth and partial stock rankings as economic objects. Second, it provides an exact certified-versus-uncertified decomposition of rank-based positions and a no-trade interpretation. Third, it brings dependence-robust decision-risk certification to cross-sectional portfolio choice without presenting the statistical method as the finance contribution. The intended finding is not that a new algorithm forecasts returns better. It is whether machine learning identifies more apparent cross-sectional differentiation than investors can reliably and economically use.
